% Source: Solutions/Kanishk-PS2-Q6-8.tex (LaTeX source).
\submittedby{Kanishk}

Suppose first that $\{x\}$ is open in the order topology. Then there exists a basic open set
$B$ such that
\[
  x\in B\subseteq\{x\}.
\]
Hence, $B=\{x\}$.

If $x$ is neither the minimum nor the maximum element, then $B$ must be an open interval
$(a,b)$, where $a<x<b$. Since $(a,b)=\{x\}$, there is no element between $a$ and $x$, or
between $x$ and $b$. Thus, $a$ is the immediate predecessor and $b$ is the immediate
successor of $x$.

If $x$ is the minimum element, then
\[
  (-\infty,b)=\{x\},
\]
so $b$ is the immediate successor of $x$. Similarly, if $x$ is the maximum element, then
\[
  (a,\infty)=\{x\},
\]
so $a$ is the immediate predecessor of $x$.

Conversely, suppose that $x$ has an immediate predecessor $a$ and an immediate successor $b$.
Then
\[
  (a,b)=\{x\}.
\]
Since $(a,b)$ is open, $\{x\}$ is open.

If $x$ is the minimum element and has an immediate successor $b$, then
\[
  (-\infty,b)=\{x\},
\]
which is open. Similarly, if $x$ is the maximum element and has an immediate predecessor $a$,
then
\[
  (a,\infty)=\{x\},
\]
which is open.

Therefore, the given conditions are necessary and sufficient for $\{x\}$ to be open.
